\documentclass[12pt,a4paper]{ctexart}

\usepackage[a4paper,margin=2.5cm]{geometry}
\usepackage{amsmath}
\usepackage{amssymb}
\usepackage{booktabs}
\usepackage{array}
\usepackage{enumitem}
\usepackage{listings}
\usepackage[dvipsnames]{xcolor}
\usepackage{hyperref}

\definecolor{codebg}{HTML}{F7F8FA}
\definecolor{codeframe}{HTML}{D9DEE7}
\definecolor{codekeyword}{HTML}{0066CC}
\definecolor{codestring}{HTML}{A31515}
\definecolor{codecomment}{HTML}{008000}

\lstdefinestyle{reportcode}{
  backgroundcolor=\color{codebg},
  basicstyle=\ttfamily\small,
  frame=single,
  rulecolor=\color{codeframe},
  columns=fullflexible,
  breaklines=true,
  breakatwhitespace=true,
  tabsize=4,
  showstringspaces=false,
  keepspaces=true,
  keywordstyle=\color{codekeyword}\bfseries,
  commentstyle=\color{codecomment},
  stringstyle=\color{codestring},
  numbers=left,
  numberstyle=\tiny\color{gray},
  numbersep=8pt,
  xleftmargin=1.8em,
  framexleftmargin=1.2em
}

\lstset{style=reportcode}
\hypersetup{
  colorlinks=true,
  linkcolor=MidnightBlue,
  urlcolor=RoyalBlue,
  citecolor=MidnightBlue
}

\title{并行程序设计与算法实验五报告\\基于 OpenMP 与 Pthreads 的并行矩阵乘法}
\author{姓名：元朗曦\quad 学号：23336294}
\date{\today}

\begin{document}

\maketitle

\section{实验目的}
本实验围绕矩阵乘法这一典型计算密集型问题，分别使用 OpenMP 与 Pthreads 进行并行实现，并分析线程数量、矩阵规模和调度策略对程序性能的影响。实验目的如下：
\begin{enumerate}[itemsep=0.3em]
  \item 掌握 OpenMP 并行循环的基本使用方法，理解 \texttt{parallel for} 与 \texttt{schedule} 子句的作用。
  \item 实现基于 Pthreads 的 \texttt{parallel\_for} 抽象，模拟 OpenMP 对并行循环的分解、分配与执行机制。
  \item 使用矩阵乘法验证并行程序的正确性，掌握通过串行程序进行结果校验的方法。
  \item 通过实验数据比较不同线程数、不同调度策略下的运行时间和 GFLOPS，分析并行加速效果。
\end{enumerate}

\section{实验原理}

\subsection*{矩阵乘法}
设矩阵 $A\in\mathbb{R}^{m\times n}$，矩阵 $B\in\mathbb{R}^{n\times k}$，矩阵乘法结果为 $C\in\mathbb{R}^{m\times k}$，则：
\[
C_{ij}=\sum_{p=0}^{n-1}A_{ip}B_{pj}.
\]

矩阵 $C$ 的每一行可以独立计算，因此可以将外层行循环分配给多个线程。不同线程写入 $C$ 的不同行，不存在写冲突，不需要加锁。

矩阵乘法的浮点运算量近似为：
\[
\text{FLOPs}=2mnk,
\]
其中一次乘法和一次加法合计为 2 次浮点运算。程序性能使用 GFLOPS 衡量：
\[
\text{GFLOPS}=\frac{2mnk}{t\times 10^9},
\]
其中 $t$ 为运行时间，单位为秒。

\subsection*{OpenMP 并行循环}
OpenMP 是共享内存并行编程接口，使用编译指导语句描述并行区域。在本实验中，OpenMP 版本使用：
\begin{lstlisting}[language=C++]
#pragma omp parallel for schedule(static)
for (int i = 0; i < m; ++i) {
    // compute row i of C
}
\end{lstlisting}

调度策略包括：
\begin{itemize}[itemsep=0.2em]
  \item \textbf{默认调度}：由 OpenMP 运行时和编译器决定循环分配方式。
  \item \textbf{静态调度}：在循环开始前将迭代均匀划分给各线程，调度开销小。
  \item \textbf{动态调度}：运行时动态分配迭代，负载均衡能力更强，但同步开销更高。
\end{itemize}

\subsection*{Pthreads parallel\_for}
Pthreads 是较底层的线程库，需要显式创建、传参和回收线程。为了模拟 OpenMP 的并行循环，本实验实现了如下形式的 \texttt{parallel\_for}：
\begin{lstlisting}[language=C++]
int parallel_for(int start, int end, int inc,
                 void* (*functor)(int, void*),
                 void* arg,
                 int num_threads,
                 Schedule schedule);
\end{lstlisting}

其中 \texttt{start}、\texttt{end} 和 \texttt{inc} 描述循环范围，\texttt{functor} 表示每次迭代执行的函数，\texttt{arg} 传递共享参数，\texttt{num\_threads} 表示期望线程数。静态调度将迭代区间平均分配给各线程；动态调度使用原子计数器为线程分配下一次迭代。

\section{实验内容}

本实验代码位于 \texttt{lab/05/src}，使用 CMake 作为构建工具。主要程序如下：
\begin{enumerate}[itemsep=0.3em]
  \item \texttt{openmp\_matmul.cpp}：基于 OpenMP 的并行矩阵乘法。
  \item \texttt{pthread\_parallel\_for\_matmul.cpp}：基于 Pthreads 的 \texttt{parallel\_for}，并以矩阵乘法作为测试样例。
  \item \texttt{benchmark.py}：随机生成矩阵规模并批量运行不同实现，用于性能测试。
\end{enumerate}

\subsection*{OpenMP 程序设计}
OpenMP 程序从命令行读取矩阵规模、线程数和调度策略，随机生成矩阵 $A$ 和 $B$，并计算矩阵 $C$。程序支持如下参数：
\begin{lstlisting}[language=bash]
./openmp_matmul <m> <n> <k> <threads> [default|static|dynamic] [--verify] [--print]
\end{lstlisting}

其中矩阵规模 $m,n,k$ 的范围为 $[128,2048]$，线程数量范围为 $[1,16]$。\texttt{--verify} 选项会执行串行矩阵乘法并比较最大绝对误差；\texttt{--print} 选项用于输出矩阵 $A$、$B$、$C$。

\subsection*{Pthreads 程序设计}
Pthreads 程序首先实现通用的 \texttt{parallel\_for}。
随后将矩阵乘法的行计算封装为 \texttt{functor}：
\begin{lstlisting}[language=C++]
void* matmul_row_functor(int row, void* raw_args) {
    // compute one row of C
    return nullptr;
}

parallel_for(0, m, 1, matmul_row_functor, &args, threads, schedule);
\end{lstlisting}

该设计将循环调度逻辑与具体计算逻辑分离，使 \texttt{parallel\_for} 可以复用于其他按循环迭代划分的任务。

\subsection*{Benchmark 流程}
实验 benchmark 按如下流程执行：
\begin{enumerate}[itemsep=0.3em]
  \item 在 $[128,2048]$ 范围内随机生成整数 $m,n,k$；
  \item 程序根据 $m,n,k$ 随机生成矩阵 $A\in\mathbb{R}^{m\times n}$ 和 $B\in\mathbb{R}^{n\times k}$；
  \item 调用并行矩阵乘法计算 $C=A\times B$；
  \item 记录运行时间和 GFLOPS；
  \item 重复选取不同随机规模，分析矩阵规模变化对性能的影响。
\end{enumerate}

benchmark 脚本支持选择实现、调度策略、样本数量和线程数。例如：
\begin{lstlisting}[language=bash]
python3 lab/05/benchmark.py \
  --samples 4 \
  --program both \
  --schedules static \
  --threads 4 \
  --csv report/random_benchmark.csv
\end{lstlisting}

\subsection*{构建与运行}
构建命令如下：
\begin{lstlisting}[language=bash]
cmake -S lab/05 -B lab/05/build -DCMAKE_BUILD_TYPE=Release
cmake --build lab/05/build -j
\end{lstlisting}

单次运行示例：
\begin{lstlisting}[language=bash]
./lab/05/build/openmp_matmul 512 512 512 4 static --verify
./lab/05/build/pthread_parallel_for_matmul 512 512 512 4 static --verify
\end{lstlisting}

\section{实验结果}

\subsection*{正确性验证}
OpenMP 与 Pthreads 两个程序均通过串行矩阵乘法校验。以 $128\times128\times128$ 矩阵为例：
\begin{lstlisting}
Task: OpenMP matrix multiplication
A: 128 x 128
B: 128 x 128
C: 128 x 128
Threads: 4
Schedule: dynamic
Time (sec): 0.000432
GFLOPS: 9.699
Verify: PASS, max_abs_diff=0.0000000000

Task: Pthreads parallel_for matrix multiplication
A: 128 x 128
B: 128 x 128
C: 128 x 128
Requested threads: 4
Actual threads: 4
Schedule: dynamic
Time (sec): 0.000618
GFLOPS: 6.791
Verify: PASS, max_abs_diff=0.0000000000
\end{lstlisting}

\subsection*{随机规模 benchmark}
本节按照 benchmark 流程随机生成多组 $m,n,k$。每组样本中，程序都会随机生成 $A$ 和 $B$，再计算 $C=A\times B$。测试线程数为 4，调度方式为 static。

\begin{table}[h]
\centering
\caption{随机矩阵规模 benchmark 结果}
\begin{tabular}{ccccccc}
\toprule
样本 & 实现 & $m$ & $n$ & $k$ & 时间（s） & GFLOPS \\
\midrule
1 & OpenMP & 458 & 179 & 1371 & 0.004380 & 51.325 \\
1 & Pthreads & 458 & 179 & 1371 & 0.004517 & 49.772 \\
2 & OpenMP & 1125 & 596 & 1402 & 0.035874 & 52.409 \\
2 & Pthreads & 1125 & 596 & 1402 & 0.034081 & 55.166 \\
3 & OpenMP & 1145 & 1441 & 1063 & 0.068255 & 51.393 \\
3 & Pthreads & 1145 & 1441 & 1063 & 0.068131 & 51.486 \\
4 & OpenMP & 359 & 197 & 420 & 0.001364 & 43.565 \\
4 & Pthreads & 359 & 197 & 420 & 0.001233 & 48.185 \\
\bottomrule
\end{tabular}
\end{table}

从结果可以看出，随机规模下两个实现均能保持较稳定的吞吐率。较小样本（如样本 4）的 GFLOPS 低于大样本，主要原因是线程创建、调度和计时开销在小规模计算中占比更高；较大样本中计算量更充分，并行开销被摊薄，因此 GFLOPS 更稳定。OpenMP 与 Pthreads 在同一随机规模下的结果接近，说明手动实现的 \texttt{parallel\_for} 能正确完成循环分解与线程执行。

\subsection*{实验结论}
\begin{enumerate}[itemsep=0.3em]
  \item OpenMP 能以较少代码实现并行矩阵乘法，适合规则循环的快速并行化。
  \item 对矩阵乘法这类每次迭代工作量接近的任务，静态调度通常具有较低开销和较好性能。
  \item Pthreads 可以构造类似 OpenMP 的 \texttt{parallel\_for} 抽象，但需要显式处理线程创建、任务划分和同步。
  \item 随着线程数增加，程序总体性能提升明显，但实际加速比会受到线程管理开销、缓存行为和硬件核心数限制。
\end{enumerate}

\end{document}
